\chapter{The Linear Algebra of the Free Block}
\label{ch:freeblock}

Every method in this book solves the same linear system many times, with the
index set changing by one between consecutive solves. This chapter assembles the
linear algebra: the saddle-point form of the working-set subproblem, the two ways
to eliminate its multipliers, the Woodbury identity, and---the part that turns an
$O(n^4)$ algorithm into an $O(n^3)$ one---how to \emph{update} a factorisation
when one index joins or leaves.

Nothing here is new mathematics. It is assembled in one place because the
algorithms of Parts~II and~III are otherwise unreadable, and because the choice
between the alternatives below is where implementations differ.

\section{The saddle-point system}

Recall the subproblem~\eqref{eq:sub}. Writing its stationarity and feasibility
conditions together gives the \emph{saddle-point} or \emph{bordered} system
\index{saddle-point system}
\begin{equation}
  \begin{bmatrix} G & -A \\ A\T & 0 \end{bmatrix}
  \begin{bmatrix} x \\ u \end{bmatrix}
  =
  \begin{bmatrix} a \\ b_\Active \end{bmatrix},
  \qquad A = C_{\cdot\Active} \in \R^{n\times k}.
  \label{eq:saddle}
\end{equation}
The coefficient matrix is symmetric but \emph{indefinite}: it has $n$ positive and
$k$ negative eigenvalues when $G \succ 0$ and $A$ has full column rank. That
indefiniteness is why one does not simply run Cholesky on it, and why the two
elimination strategies below exist.

\begin{proposition}[Nonsingularity]
\label{prop:saddlenonsing}
If $G \succ 0$ and $A$ has full column rank then the matrix in~\eqref{eq:saddle}
is nonsingular.
\end{proposition}

\begin{proof}
Suppose $(v,w)$ lies in the kernel: $Gv = Aw$ and $A\T v = 0$. Then
$v\T G v = v\T A w = (A\T v)\T w = 0$, so $v = 0$ by definiteness of $G$, whence
$Aw = 0$ and $w = 0$ by full column rank.
\end{proof}

The same three-line argument recurs in Part~III for the bordered system of the
parametric theory (Lemma~\ref{lem:affine}), where $G$ is replaced by the reduced
Hessian on the free block. It is worth learning once.

\section{Two eliminations}

\subsection{Eliminate $x$: the dual Hessian route}

Solve the first block row for $x$ and substitute:
\[
  x = G^{-1}(a + Au)
  \quad\Longrightarrow\quad
  \underbrace{(A\T G^{-1} A)}_{k \times k}\,u = b_\Active - A\T G^{-1}a .
\]
This is Proposition~\ref{prop:sub} again. Cost: one Cholesky of $G$ ($O(n^3)$,
once), $k$ triangular solves to form $G^{-1}A$ ($O(n^2 k)$), and a $k\times k$
Cholesky ($O(k^3)$). The route is natural when $k$ is small and $G$ is factored
once and reused across many working sets---which is exactly the situation of
Chapters~\ref{ch:guessing} and~\ref{ch:warmstart}.

\subsection{Eliminate $u$: the reduced Hessian route}

Alternatively, let the columns of $Z \in \R^{n \times (n-k)}$ be a basis of
$\nullsp(A\T)$, write $x = x_p + Z y$ for any particular $x_p$ with
$A\T x_p = b_\Active$, and substitute into the objective. The result is an
unconstrained problem in $y$ with Hessian $Z\T G Z$, the \emph{reduced Hessian},
positive definite because $G$ is. Cost: forming $Z$, then an $(n-k)$-dimensional
solve.
\index{reduced Hessian}

Which is better depends on $k$ versus $n - k$, and the answer flips between the
two regimes this book cares about. A general QP with few active constraints has
$k \ll n$: eliminate $x$. A long-only portfolio at the optimum has $k \approx n$,
almost every bound binding, so $n - k$ is tiny: eliminate $u$, or equivalently
work on the free block directly, which is what Chapters~\ref{ch:krylov}
and~\ref{ch:operator} do.

The method of Chapter~\ref{ch:walking} does something cleverer than either: it
carries a factorisation from which \emph{both} the reduced inverse Hessian and the
dual Hessian's Cholesky factor can be read off, and updates it in $O(n^2)$ per
step. Section~\ref{sec:updating} below prepares that.

\section{Block elimination in the parametric setting}

Part~III needs a variant worth recording now. When the right-hand side splits into
a constant part and a $\lambda$-linear part, the system is solved for \emph{two}
right-hand sides at once, against one factorisation. With
\[
  Y = G_{\Free\Free}^{-1} C_\Free\T, \qquad
  z = G_{\Free\Free}^{-1} r_1,
\]
obtained from a single multi-column solve, the multipliers follow from the
$k \times k$ Schur complement $S = C_\Free\,G_{\Free\Free}^{-1} C_\Free\T = C_\Free Y$:
\begin{equation}
  \nu = S^{-1}\bigl(C_\Free z - r_2\bigr),
  \qquad
  x_\Free = z - Y\nu .
  \label{eq:blockelim}
\end{equation}
\index{Schur complement}
The dominant cost is the multi-right-hand-side solve against the free block; the
Schur solve is $k\times k$, and for the single budget constraint of a portfolio
problem $k = 1$, so it is a scalar division. This is the inner loop of the
Critical Line Algorithm (Chapter~\ref{ch:parametric}) and of the equality-augmented
solver of Chapter~\ref{ch:krylov}, in both cases with the same structure: the
multipliers are recovered \emph{analytically} and never appear as unknowns in the
linear solve, so every system actually solved stays symmetric positive definite.

That last sentence is the point of the manoeuvre. An indefinite system needs
$LDL\T$ with pivoting; an SPD system takes Cholesky, or conjugate gradients, or a
Woodbury formula. Keeping the solves SPD is what lets Chapter~\ref{ch:krylov} use
a Krylov method at all.

\section{The Woodbury identity}

\begin{proposition}[Sherman--Morrison--Woodbury]
\label{prop:woodbury}
\index{Woodbury identity}
For invertible $D \in \R^{n\times n}$ and $\Delta \in \R^{k\times k}$ and
$U \in \R^{n\times k}$,
\[
  (D + U\Delta U\T)^{-1}
   = D^{-1} - D^{-1}U\bigl(\Delta^{-1} + U\T D^{-1} U\bigr)^{-1} U\T D^{-1},
\]
whenever the inverses on the right exist.
\end{proposition}

\begin{proof}
Multiply the right-hand side by $D + U\Delta U\T$ and simplify, using
$\Delta U\T D^{-1}U = (\Delta^{-1} + U\T D^{-1}U)\Delta - I$ to collapse the cross
terms.
\end{proof}

The identity converts an $n\times n$ inverse into a $k\times k$ one, and it is the
engine of Chapter~\ref{ch:operator}. Its importance here is that
\emph{diagonal-plus-low-rank} is the form that structured models actually take: a
factor risk model $\Sig = \diag(d) + U\Lambda U\T$ with $k \ll n$ factors; a
random-matrix-cleaned correlation $\bar\mu I + U_k\Delta_k U_k\T$; a Nyström kernel
approximation; a spiked covariance model. All four are the same object, and all
four invert in $O(nk^2 + k^3)$ instead of $O(n^3)$.

One subtlety that catches implementers. When the identity is applied to a
\emph{restricted} block $\Active$, the restricted factor $U_\Active = U[\Active,:]$
is generally \emph{not} orthonormal even if $U$ was:
$U_\Active\T U_\Active \ne I_k$. The correction matrix
$W = \Delta^{-1} + U_\Active\T D_\Active^{-1} U_\Active$ accounts for this, and it
is what must be formed and factored per solve. Its cost, $O(n_a k^2 + k^3)$, is
what Proposition~\ref{prop:crossover} in Chapter~\ref{ch:operator} weighs against
assembling and factoring the $n_a \times n_a$ block directly.

\section{Updating a factorisation}
\label{sec:updating}

An active-set method changes the working set by one index per step. Recomputing a
factorisation from scratch costs $O(nk^2)$; \emph{updating} it costs $O(n^2)$ or
$O(nk)$. Over the $O(n)$ steps of a run that is the difference between $O(n^4)$
and $O(n^3)$, and it is the entire reason factorisations are carried rather than
rebuilt.

Both updates below work the same way. Suppose we carry a matrix $J$ satisfying
$J\T G J = I$ together with an upper triangular $R$ with $J\T A = [R;0]$
(Chapter~\ref{ch:walking} explains why these are the right invariants). The first
condition is preserved by post-multiplying $J$ with \emph{any} orthogonal $W$,
since $(JW)\T G (JW) = W\T W = I$. So the update's only task is to choose $W$ so
that the triangular shape of the second condition is restored. Because
$J' = JW$ gives ${J'}\T A = W\T(J\T A)$, the same $W\T$ acts on the rows of $R$.

\subsection{Insertion}
\index{factorisation update!insertion}

Append a column $n_*$ to $A$. With $d = J\T n_*$ split as $(d_1,d_2)$ at index $k$,
\[
  J\T [\,A \ \ n_*\,] = \begin{bmatrix} R & d_1 \\ 0 & d_2 \end{bmatrix},
\]
which fails to be triangular only in the trailing $n-k-1$ entries of $d_2$. Choose
orthogonal $W_2$ with $W_2\T d_2 = \alpha e_1$ and set $W = \diag(I_k, W_2)$. Then
the new $R$ is $R$ bordered by $d_1$ and $\alpha$, and the leading $k$ columns of
$J$ are untouched.

Two implementations of $W_2$ compete. A \emph{chain of Givens rotations}, one per
trailing component, is $n-k-1$ operations of $O(n)$ each. A single
\emph{Householder reflection}
\begin{equation}
  W_2 = I - \frac{2}{\beta}\, v v\T,
  \qquad v = d_2 - \alpha e_1,
  \qquad \beta = v\T v,
  \qquad \alpha = \pm\norm{d_2},
  \label{eq:householder}
\end{equation}
achieves the same reduction in two operations of $O(n(n-k))$: a matrix--vector
product and a rank-one update. The arithmetic is comparable; the number of
operations is not, and on modern hardware the operation count is what one pays.

\begin{remark}[The sign of $\alpha$]
\label{rem:sign}
Numerical analysis prescribes $\alpha = -\sign((d_2)_1)\norm{d_2}$, making
$v_1 = (d_2)_1 - \alpha$ a sum of like-signed terms. A code that takes the opposite
sign---to match a Givens chain's sign convention, say---walks into the very
cancellation the prescription exists to avoid, severely when $d_2$ is already close
to a positive multiple of $e_1$. The cure is algebraic, not a change of convention:
with $\tau = \norm{(d_2)_{2:}}^2$ and $\nu = \norm{d_2}$,
\begin{equation}
  v_1 = (d_2)_1 - \sign((d_2)_1)\,\nu
      = \frac{(d_2)_1^2 - \nu^2}{(d_2)_1 + \sign((d_2)_1)\nu}
      = \frac{-\sign((d_2)_1)\,\tau}{|(d_2)_1| + \nu},
  \label{eq:nocancel}
\end{equation}
in which every operation combines like-signed quantities. Chapter~\ref{ch:numerics}
returns to this pattern of rewriting a difference as a quotient.
\end{remark}

\subsection{Deletion}
\index{factorisation update!deletion}

Delete column $p$ of $A$. This deletes column $p$ of $R$ and shifts the rest left,
leaving a matrix that is upper triangular except for one subdiagonal, with
offending entries $\tilde R_{j+1,j}$ for $j = p,\dots,k-1$. These lie in
\emph{distinct} columns, and no single reflection annihilates components of
distinct vectors. What restores the shape is a \emph{chase}: for $j = p,\dots,k-1$
in order, a $2\times2$ orthogonal acting on rows $j,j+1$ annihilates
$\tilde R_{j+1,j}$, and the same transformation is applied to columns $j, j+1$
of $J$.

A useful implementation trick: take the $2\times 2$ block to be the \emph{symmetric}
reflection
\begin{equation}
  W_j^{(2)} = \begin{bmatrix} c & s \\ s & -c \end{bmatrix},
  \qquad c = \frac{x}{h},\quad s = \frac{y}{h},
  \qquad h = \sign(x)\sqrt{x^2+y^2},
  \label{eq:refl2}
\end{equation}
rather than a Givens rotation. The update applies $W$ to the columns of $J$ and
$W\T$ to the rows of $R$; symmetry makes those the same matrix, so one $2\times2$
serves both sides. A Givens rotation is not symmetric and needs its transpose on
the $R$ side. The two differ by a sign, and Proposition~\ref{prop:invariance} in
Chapter~\ref{ch:walking} shows that difference is free.

\begin{remark}[Why the reflections cannot be blocked]
\label{rem:wy}
\index{WY representation}
A sequence of Householder reflections is normally applied as a block: the WY
representation \cite{bischof1987,schreiber1989} writes $W_1\cdots W_j$ as
$I + WY\T$, so $j$ reflections reach a matrix in one level-3 update rather than
$j$ level-2 ones. That is why blocked QR is fast, and it is the first thing to
reach for here.

It is not available, for a structural reason. Blocking requires the $j$
reflections to be known before any is applied. In an active-set iteration the
reflection that inserts a constraint is determined by $d = J\T n_*$, which depends
on $J$ \emph{after} the previous insertion; and between two insertions the solver
reads quantities off the updated factors to decide which constraint enters next,
and whether it enters at all. The updates are interleaved with reads of what they
produce, so there is no window in which two reflections are simultaneously known
and unapplied. The same obstruction governs the deletion chase, whose
transformations read their parameters from entries the previous one wrote.

\emph{The method therefore stays a level-2 computation no matter how it is coded.}
This is a real ceiling, and it is one reason Chapter~\ref{ch:measurement} argues
that iteration counts, not flop counts, are the currency at these sizes.
\end{remark}

\section{Conditioning, briefly}

The spectral condition number $\kappa(G) = \lambda_{\max}(G)/\lambda_{\min}(G)$
governs both the accuracy of a direct solve and the iteration count of an
iterative one. Two facts are used repeatedly later.

\begin{proposition}[Cauchy interlacing, the consequence we need]
\label{prop:interlace}
If $G \succ 0$ then $\lambda_{\min}(G_{\Free\Free}) \ge \lambda_{\min}(G)$ and
$\lambda_{\max}(G_{\Free\Free}) \le \lambda_{\max}(G)$ for every index set $\Free$.
Hence $\kappa(G_\Free) \le \kappa(G)$.
\end{proposition}

\begin{proof}
Both bounds follow from the Rayleigh quotient characterisation restricted to
vectors supported on $\Free$.
\end{proof}

So a free-block solve is never worse conditioned than the full problem---a small
mercy that Chapter~\ref{ch:krylov} leans on when it bounds an inner CG iteration
count uniformly over the working sets a loop might visit.

\begin{proposition}[Weyl, the consequence we need]
\label{prop:weyl}
\index{Weyl's inequality}
For $\alpha \in (0,1)$ and $R\T R \succ 0$, the convex combination
$G_\alpha = (1-\alpha)G + \alpha R\T R$ satisfies
\[
  \kappa(G_\alpha) \;\le\;
  \frac{(1-\alpha)\lambda_{\max}(G) + \alpha\,\lambda_{\max}(R\T R)}
       {\alpha\,\lambda_{\min}(R\T R)},
\]
which decreases monotonically in $\alpha$. For $R\T R = \bar\lambda I$ the bound is
$\tfrac{1-\alpha}{\alpha}\,\lambda_{\max}(G)/\bar\lambda + 1$.
\end{proposition}

\begin{proof}
Weyl's inequality \cite{hornjohnson2013} bounds $\lambda_{\max}(G_\alpha)$ above by
the sum of the two terms' maxima and $\lambda_{\min}(G_\alpha)$ below by
$\alpha\lambda_{\min}(R\T R)$, the first term contributing at least zero.
\end{proof}

This is the whole mathematical content of ``regularisation improves conditioning'',
and Chapter~\ref{ch:conditioning} spends it three times over: on Tikhonov
regularisation, on Ledoit--Wolf shrinkage, and on the observation that a
statistically motivated shrinkage intensity conditions the solver as a free
by-product.

\begin{notes}
Golub and Van Loan \cite{golub2013} covers saddle-point systems, Householder and
Givens transformations, and factorisation updating; Benzi, Golub and Liesen
\cite{benzi2005} is the survey of saddle-point problems specifically. Higham
\cite{higham2002} is the reference for the cancellation analysis behind
Remark~\ref{rem:sign}. The WY representation is \cite{bischof1987,schreiber1989};
the argument in Remark~\ref{rem:wy} that it is structurally unavailable here is
from the Goldfarb--Idnani note \cite{schmelzer2026quadprog}.
\end{notes}

\begin{exercises}

\begin{exercise}
Show that the saddle-point matrix in~\eqref{eq:saddle} has exactly $n$ positive and
$k$ negative eigenvalues when $G \succ 0$ and $A$ has full column rank. (Hint: use
the block $LDL\T$ factorisation with the Schur complement $-A\T G^{-1}A$.)
\end{exercise}

\begin{exercise}
Count the flops of the two eliminations of Section~2 for given $n$ and $k$,
assuming $G$ is already factored. At what $k/n$ do they cross?
\end{exercise}

\begin{exercise}
Prove Proposition~\ref{prop:woodbury} by verifying the product both ways. Then
specialise to $k = 1$ to recover the Sherman--Morrison formula.
\end{exercise}

\begin{exercise}
Let $U \in \R^{n\times k}$ have orthonormal columns and let $\Active$ index $n_a$
rows. Show $\norm{U_\Active\T U_\Active}_2 \le 1$, with equality iff some unit
vector in $\range(U\T)$ is supported on $\Active$. Explain why the correction
matrix $W$ of Section~4 is nevertheless always invertible when $\Delta \succ 0$.
\end{exercise}

\begin{exercise}
Implement the Householder insertion~\eqref{eq:householder} both with the naive
$v_1 = (d_2)_1 - \sign((d_2)_1)\norm{d_2}$ and with the rewriting~\eqref{eq:nocancel}.
Construct $d_2$ close to a positive multiple of $e_1$ and measure the relative
error in $\norm{W_2\T d_2 - \alpha e_1}$ for both.
\end{exercise}

\begin{exercise}
Show that the deletion chase of Section~5.2 requires exactly $k - p$ plane
transformations, and that each is determined by entries written by the previous
one. Conclude that no reordering makes it parallel.
\end{exercise}

\begin{exercise}[Implementation]
Write a routine that maintains a Cholesky factor of $A\T G^{-1}A$ under insertion
and deletion of columns of $A$, and a second routine that recomputes it from
scratch. On random data with $n = 200$, insert $100$ columns one at a time and
compare total runtime and final accuracy (measured against a fresh factorisation).
At what point does the accumulated error in the updated factor become visible?
\end{exercise}

\end{exercises}
